\begin{exercice}\label{exo4-4} (\minsyndical)

\begin{enumerate}

\item Soit $X : \Omega \rightarrow \mathbb{R}$ une variable al\'eatoire constante, i.e. pour tout $\omega \in \Omega$, $X(\omega) = c$, avec $c\in\mathbb{R}$. Montrer que $E[X] = c$, $Var(X) = 0$, $\sigma (X) = 0$.

\item Soit $X : \Omega \rightarrow \mathbb{R}$ une variable al\'eatoire discr\`ete, $a$ et $b$ deux r\'eels, et $\phi,\psi$ deux fonctions de $\mathbb{R}$ dans $\mathbb{R}$. Montrer que :
\[
E [ a \phi(X) + b \psi(X) ] =
a E [\phi(X)] + b  E [ \psi(X) ].
\]

\item Montrer que la variance peut se calculer comme suit :
\[
Var (X) = E[X^2] - E[X]^2.
\]

\end{enumerate}

%\corrref{TD1_1}

\end{exercice}
